% \subsubsection{Discussion}

\red{[TODO: discussion rewritten from the new results — learning pays against most opponents, loses against another learner (the sign reversal), and the price of safety. The old claims that FP ``consistently outperforms'' all baselines, wins on fuel, and that random beats greedy ``at every horizon'' were removed: the first is contradicted by the sign reversal, the others were never re-measured.]}

The advantage is primarily informational: all players share the same action space and $\Delta V$ constraints, yet FP players achieve better outcomes by learning the opponent's behavior online.

In future work, the fictitious-play framework can be extended with Bayesian meta-strategy reasoning, in which the agent maintains a posterior belief over a library of opponent types and best-responds by marginalizing over this belief. This would enable the agent to distinguish not only \textit{which} vertices the opponent favors but \textit{why}, and accelerate learning against dynamic strategies.